1.1 Non‑Dual Algebra

`tex
\subsection{Non-Dual Algebra}

\begin{definition}[Non-Dual Algebra]
A non-dual algebra is a triple $(\mathcal{N}, \circ, \mathcal{K})$ where:
\begin{itemize}
  \item $\mathcal{N}$ is a set of operators,
  \item $\circ : \mathcal{N} \times \mathcal{N} \to \mathcal{N}$ is an associative composition,
  \item $\mathcal{K}$ is a kernel that is not required to be preserved under composition.
\end{itemize}
Duals $U^\*$ need not exist.
\end{definition}

\begin{definition}[Non-Dual Curvature]
For $U,V \in \mathcal{N}$, the curvature is the commutator
\[
K(U,V) = U \circ V - V \circ U.
\]
\end{definition}

\begin{axiom}[Non-Dual Drift]
Non-dual curvature induces drift:
\[
K(U,V) \neq 0 \Rightarrow \text{drift accumulates along operator chains}.
\]
\end{axiom}

\begin{axiom}[Kernel Non-Preservation]
For $p \in \mathcal{K}$,
\[
U(p) \notin \mathcal{K}
\]
is permitted.
\end{axiom}
`

---

1.2 Collapse Algebra

`tex
\subsection{Collapse Algebra}

\begin{definition}[Collapse Operator]
A collapse operator is a map $C : X \to X$ such that
\[
C(x) = x' \quad \text{with} \quad x' \le x
\]
under a partial order representing structural contraction.
\end{definition}

\begin{definition}[Collapse Residual]
The residual of collapse is
\[
R_C(x) = x - C(x).
\]
\end{definition}

\begin{axiom}[Collapse Monotonicity]
For all $x,y$,
\[
x \le y \Rightarrow C(x) \le C(y).
\]
\end{axiom}

\begin{axiom}[Residual Boundedness]
Collapse residuals satisfy
\[
R_C(x) \le x.
\]
\end{axiom}
`

---

1.3 Holonomy Algebra

`tex
\subsection{Holonomy Algebra}

\begin{definition}[Holonomy]
Given operators $U,V,W$, holonomy is the loop
\[
H = U^{-1} \circ V \circ U \circ W^{-1},
\]
where inverses may be partial or absent.
\end{definition}

\begin{definition}[Holonomy Curvature]
Holonomy curvature is defined as
\[
\Omega = H - I,
\]
where $I$ is the identity operator (if it exists).
\end{definition}

\begin{axiom}[Holonomy Consistency]
Holonomy curvature satisfies
\[
\Omega(U \circ V) = \Omega(U) + \Omega(V) + K(U,V).
\]
\end{axiom}

\begin{axiom}[Holonomy Drift]
Non-zero curvature induces drift along loops.
\end{axiom}
`

---

1.4 Modal‑Truth Algebra

`tex
\subsection{Modal-Truth Algebra}

\begin{definition}[Truth Lattice]
A truth lattice $(\mathcal{T}, \wedge, \vee)$ is a bounded lattice with meet $\wedge$
and join $\vee$.
\end{definition}

\begin{definition}[Truth Residue]
Truth residue is defined as
\[
R_T(t) = t \wedge \neg t',
\]
for some $t'$ induced by phase geometry.
\end{definition}

\begin{definition}[Modal Lift]
A modal lift is a map
\[
L : \mathcal{T} \to \mathcal{M}
\]
into a modal frame $\mathcal{M}$.
\end{definition}

\begin{axiom}[Lattice Closure]
Truth lattices satisfy
\[
t1, t2 \in \mathcal{T} \Rightarrow t1 \wedge t2, t1 \vee t2 \in \mathcal{T}.
\]
\end{axiom}

\begin{axiom}[Modal Invariance]
Modal lifts preserve truth ordering:
\[
t1 \le t2 \Rightarrow L(t1) \le L(t2).
\]
\end{axiom}
`

---

1.5 Spectral‑Bundle Algebra

`tex
\subsection{Spectral-Bundle Algebra}

\begin{definition}[Bundle]
A bundle is a triple $(B, \pi, F)$ where:
\begin{itemize}
  \item $B$ is the total space,
  \item $\pi : B \to X$ is a projection,
  \item $F$ is the fiber.
\end{itemize}
\end{definition}

\begin{definition}[Spectral Lift]
A spectral lift is a map
\[
S : X \to B
\]
such that $\pi(S(x)) = x$.
\end{definition}

\begin{definition}[Spectral Curvature]
Spectral curvature is defined as
\[
\kappa_S(x) = \| S(x) - S'(x) \|,
\]
for two lifts $S,S'$.
\end{definition}

\begin{axiom}[Spectral Boundedness]
Spectral curvature satisfies
\[
\kappa_S(x) < \infty.
\]
\end{axiom}

\begin{axiom}[Projection Consistency]
For all $x$,
\[
\pi(S(x)) = x.
\]
\end{axiom}
`

---

